\documentclass[10pt,a4paper]{amsart}
\usepackage{amsmath,amssymb,amsthm,graphicx}
\usepackage{verbatim}
\newenvironment{reference}{\comment}{\endcomment}
%\newenvironment{reference}{}{}
\newenvironment{slogan}{\comment}{\endcomment}
\newenvironment{history}{\comment}{\endcomment}

% For commutative diagrams we use Xy-pic
\usepackage[all]{xy}

% We use 2cell for 2-commutative diagrams.
\xyoption{2cell}
% Theorem environments.
%
\theoremstyle{plain}
\newtheorem{theorem}[subsection]{Theorem}
\newtheorem{proposition}[subsection]{Proposition}
\newtheorem{lemma}[subsection]{Lemma}

\theoremstyle{definition}
\newtheorem{definition}[subsection]{Definition}
\newtheorem{example}[subsection]{Example}
\newtheorem{exercise}[subsection]{Exercise}
\newtheorem{situation}[subsection]{Situation}

\theoremstyle{remark}
\newtheorem{remark}[subsection]{Remark}
\newtheorem{remarks}[subsection]{Remarks}

\numberwithin{equation}{subsection}

% Macros
%
\def\lim{\mathop{\mathrm{lim}}\nolimits}
\def\colim{\mathop{\mathrm{colim}}\nolimits}
\def\Spec{\mathop{\mathrm{Spec}}}
\def\Hom{\mathop{\mathrm{Hom}}\nolimits}
\def\Ext{\mathop{\mathrm{Ext}}\nolimits}
\def\SheafHom{\mathop{\mathcal{H}\!\mathit{om}}\nolimits}
\def\SheafExt{\mathop{\mathcal{E}\!\mathit{xt}}\nolimits}
\def\Sch{\mathit{Sch}}
\def\Mor{\mathop{\mathrm{Mor}}\nolimits}
\def\Ob{\mathop{\mathrm{Ob}}\nolimits}
\def\Sh{\mathop{\mathit{Sh}}\nolimits}
\def\NL{\mathop{N\!L}\nolimits}
\def\CH{\mathop{\mathrm{CH}}\nolimits}
\def\proetale{{pro\text{-}\acute{e}tale}}
\def\etale{{\acute{e}tale}}
\def\QCoh{\mathit{QCoh}}
\def\Ker{\mathop{\mathrm{Ker}}}
\def\Im{\mathop{\mathrm{Im}}}
\def\Coker{\mathop{\mathrm{Coker}}}
\def\Coim{\mathop{\mathrm{Coim}}}

% Boxtimes
%
\DeclareMathSymbol{\boxtimes}{\mathbin}{AMSa}{"02}

%
% Macros for moduli stacks/spaces
%
\def\QCohstack{\mathcal{QC}\!\mathit{oh}}
\def\Cohstack{\mathcal{C}\!\mathit{oh}}
\def\Spacesstack{\mathcal{S}\!\mathit{paces}}
\def\Quotfunctor{\mathrm{Quot}}
\def\Hilbfunctor{\mathrm{Hilb}}
\def\Curvesstack{\mathcal{C}\!\mathit{urves}}
\def\Polarizedstack{\mathcal{P}\!\mathit{olarized}}
\def\Complexesstack{\mathcal{C}\!\mathit{omplexes}}
% \Pic is the operator that assigns to X its picard group, usage \Pic(X)
% \Picardstack_{X/B} denotes the Picard stack of X over B
% \Picardfunctor_{X/B} denotes the Picard functor of X over B
\def\Pic{\mathop{\mathrm{Pic}}\nolimits}
\def\Picardstack{\mathcal{P}\!\mathit{ic}}
\def\Picardfunctor{\mathrm{Pic}}
\def\Deformationcategory{\mathcal{D}\!\mathit{ef}}
\usepackage{cleveref}
\crefname{originalsection}{original section}{original sections}
\Crefname{originalsection}{Original section}{Original sections}
\makeatletter
\newlabel{algebra-example-localize-at-prime}{{02C5}{source}{}{sourcepointer-1}{}}
\newlabel{algebra-example-localize-at-prime@cref}{{[originalsection][0][]02C5}{[1][0][]source}{}{}{}}
\makeatother
\crefname{originalreference}{original reference}{original references}
\Crefname{originalreference}{Original reference}{Original references}
\makeatletter
\newlabel{proposition-localization-and-serre-quotients}{{0EA5}{source}{}{sourcepointer-2}{}}
\newlabel{proposition-localization-and-serre-quotients@cref}{{[originalreference][0][]0EA5}{[1][0][]source}{}{}{}}
\makeatother
\title{Mod\_A modulo torsion modules is canonically equivalent to Q(A)-modules, Q(A) the total quotient ring}
\author{The Stacks Project Authors}
\date{stacks-project:examples.tex:local-snapshot-20261012; GNU FDL 1.2 or later}
\begin{document}
\maketitle
\noindent\textit{Editorial note.} All delivered mathematical text is the author's own, copied from the pinned Stacks Project chapter file examples.tex; no statement or proof step was rewritten, reordered or removed.\par\medskip
\noindent\textit{Editorial note.} The counted proof is the author's complete argument: The proof is a two-line specialization of proposition-localization-and-serre-quotients (delivered card L48282 of this lane); the nontrivial identification 'torsion iff every element annihilated by a non-zerodivisor' is stated in the proof. Content verbatim.\par\medskip
\noindent\textit{Editorial note.} The label proposition-quotient-by-torsion-modules names the counted statement, and the proof environment immediately following it in the source is delivered verbatim as its proof.\par\medskip
\noindent\textit{Editorial note.} The counted statement is carried by the author's own proposition environment in the chapter source.\par\medskip
\noindent\textit{Allowed standard prerequisites (editorial).}\begin{itemize}
\item Serre subcategories and quotient categories; the localization equivalence Mod\_A/T -> Mod\_\{S\textasciicircum{}\{-1\}A\} for a Serre subcategory of S-torsion modules (same-chapter proposition-localization-and-serre-quotients, delivered as L48282)
\item The total quotient ring Q(A) = S\textasciicircum{}\{-1\}A with S the non-zerodivisors (Algebra, example-localize-at-prime); torsion module = every element annihilated by some non-zerodivisor
\end{itemize}\medskip
\ifdefined\hypertarget\hypertarget{sourcepointer-1}{}\fi
\noindent\textit{Source reference.} Original section 02C5: sources/stacks-chow/tags.txt, source lines 3072--3072. Prerequisite pointer invoked at 6820; it names a standard fact of a sibling chapter that the counted argument uses, and every step of the counted proof built on it is retained verbatim in the delivered body.\par\medskip
\ifdefined\hypertarget\hypertarget{sourcepointer-2}{}\fi
\noindent\textit{Source reference.} Original reference 0EA5: sources/stacks-examples/pinned-original-source.tex, source lines 6784--6784. Prerequisite pointer invoked at 6830; it names a standard fact of the same chapter that the counted argument uses, and every step of the counted proof built on it is retained verbatim in the delivered body.\par\medskip
\begin{proposition}
\label{proposition-quotient-by-torsion-modules}
Let $A$ be a ring. Let $Q(A)$ denote its total quotient ring
(as in Algebra, Example \ref{algebra-example-localize-at-prime}). Let
$\text{Mod}_A$ denote the category of $A$-modules and $\mathcal{T}$ its
Serre subcategory of torsion modules. Let $\text{Mod}_{Q(A)}$
denote the category
of $Q(A)$-modules. Then there is a canonical equivalence
$\text{Mod}_A/\mathcal{T} \rightarrow \text{Mod}_{Q(A)}$.
\end{proposition}

\begin{proof}
Follows immediately from applying Proposition
\ref{proposition-localization-and-serre-quotients} to the multiplicative subset
$S = \{f \in A \mid f \text{ is not a zerodivisor in }A\}$, since a module is a
torsion module if and only if all of its elements are each annihilated by some
element of $S$.
\end{proof}

\end{document}
